% Wave-14 Q4/20 --- Kapranov voice --- slides 723-852 extraction
% Primary source: /tmp/automorphic-slides.txt (Lorgat 2020 presentation).
% Cross-reference: Lorgat 2020 paper; Gritsenko-Clery 2008; Oberdieck-Pixton 2017.

\providecommand{\LII}{\Lambda^{2,1}_{\mathrm{II}}}
\providecommand{\Dfive}{\Delta_5}
\providecommand{\phizo}{\phi_{0,1}}
\providecommand{\kBKM}{\kappa_{\mathrm{BKM}}}
\providecommand{\kch}{\kappa_{\mathrm{ch}}}
\providecommand{\gDfive}{\mathfrak{g}_{\Dfive}}

\section*{Slides 723--852: the $\phizo$-Borcherds lift to $\Dfive$}

Lines $723$--$852$ of the automorphic-corrections presentation execute a
single argument: the Borcherds-product proof that the denominator of the
GKM superalgebra $\gDfive$ equals $\tfrac{1}{64}\Dfive(2z)$. The block
opens where $\Dfive(2z) = 64\,\Phi(z)$ is announced (Borcherds--Gritsenko--Nikulin) and
closes just before the Jacobi-group + $S$-involution generate
$\mathrm{PSp}_4(\mathbb{Z})$. The interval is the \emph{engine room} of
$N=1$ in the universal weight identity $\kBKM(\Phi_N) = c_N(0)/2$, not
the Conjecture~1 assembly floor.

\subsection*{The master denominator}

Fourier--Jacobi expansion of $\Delta_{12} = \phi_{12,1}/\delta_{12}$ through the
genus-two cusp form gives the weak Jacobi form
\[
  \phizo(z_1, z_2) \;=\; (r^{-1} + 10 + r) + q(10r^{-2} - 64 r^{-1} + 108 - 64 r + 10 r^2) + \cdots
  \qquad (q = e^{2\pi i z_1},\ r = e^{2\pi i z_2}),
\]
with Fourier coefficients $f(n, l)$, $\phizo = \sum f(n,l)\,e^{2\pi i (n z_1 + l z_2)}$.
The slide-723 theorem is the product
\[
  \tfrac{1}{64}\,\Dfive(2z) \;=\; \exp\!\big(\pi i (z_1 + z_2 + z_3)\big)
  \prod_{\substack{(n,l,m) \succ 0}} \big(1 - \exp(2\pi i (n z_1 + l z_2 + m z_3))\big)^{f(nm, l)},
\]
with the standard positivity convention: $n \geq 0$, $m \geq 0$, $l$ arbitrary
when $n > 0$ or $m > 0$; and $l < 0$ when $n = m = 0$. Matching exponents
$\operatorname{mult} \alpha = f(nm, l)$ to the Weyl--Kac--Borcherds numerator is
the root-superdimension content of $\gDfive$.

\subsection*{Convergence: Kac asymptotics}

Expressing Fourier coefficients of $\phi_{12,1}$ as a linear combination of
Jacobi theta functions pins
\[
  f(n, l) \;=\; O\!\big(\exp(\sqrt{4n - l^2})\big),
\]
which, via the Kac estimate on infinite-dimensional Lie-algebra characters,
forces convergence of the Borcherds product on a neighbourhood of the
zero-dimensional cusp of $\mathrm{Sp}_4(\mathbb{Z})$. The asymptotic is
sharp: $4nm - l^2$ is the Lorentzian norm on $\LII$, and positivity of
this form at imaginary roots is the $(\alpha, \alpha) < 0$ input to the
Borcherds character formula.

\subsection*{Factorisation $A \cdot B$}

The product splits along the cusp-decomposition of $\LII \cap \mathbb{R}_{>0}\mathcal{P}_{\mathrm{II}}$:
\[
  A \;=\; \prod_{\substack{n > 0,\ l \in \mathbb{Z} \\ \text{or } n = 0,\ l < 0}}
    \big(1 - e^{2\pi i(n z_1 + l z_2)}\big)^{f(0, l)}, \qquad
  B \;=\; \prod_{n \geq 0,\, m > 0,\, l \in \mathbb{Z}} \big(1 - e^{2\pi i(n z_1 + l z_2 + m z_3)}\big)^{f(nm, l)}.
\]
The constant-term data is $f(0, 0) = 10$, $f(0, -1) = 1$, $f(0, l) = 0$ for
$l < -1$: the $10$ is the $\kBKM(\Phi_1) = c_1(0)/2 = 5$ Borcherds weight doubled,
and the $1$ at $l = -1$ is the odd-simple-root imaginary generator that
forces $\gDfive$ to be a Lie \emph{super}algebra.

\subsection*{Hecke-minus structure on $B$}

Let $T_-(m)$ denote the $\mathrm{GL}_2(\mathbb{Z})$-minus-embedding Hecke
operator on Jacobi forms: $\phi\ \text{weight }k\ \text{index }t \mapsto (\phi \mid_k T_-(m))$
of index $mt$. The slide gives
\[
  \log B \;=\; -\sum_{m \geq 1} \frac{1}{m^2}\,\big(\phizo\, e^{2\pi i t z_3} \mid_0 T_-(m)\big),
\]
exhibiting $B$ as the exponential of a $\Gamma_\infty$-invariant sum. The
$1/m^2$ weighting is the Borcherds-lift signature at weight zero; it is the
arithmetic incarnation of the logarithmic derivative $\partial \log \Phi$
through the maximal parabolic.

\subsection*{Factor $A$: Dedekind $\eta^9$ times Jacobi theta}

Explicit closed form:
\[
  A \;=\; \psi_{5, \frac{1}{2}}\,\exp(2\pi i t z_3) \;=\; \eta(z_1)^9\,\nu_{1,1}(z_1, z_2),
\]
with $\eta(\tau) = e^{\pi i \tau / 12}\prod_{n \geq 1}(1 - e^{2\pi i n \tau})$ and
\[
  \nu_{1,1}(z_1, z_2) \;=\; \sum_{n \in \mathbb{Z}} (-1)^n\,\exp\!\Big(\tfrac{\pi i}{4}(2n+1)^2 z_1 + \pi i (2n+1) z_2\Big).
\]
Weight $\tfrac{9}{2} + \tfrac{1}{2} = 5$ matches $\Dfive$. The $\eta^9$
exponent is the $c_1(0) - 1 = 9$ combinatorics encoded in the $\tau$-invariant.

\subsection*{Modularity: $\Gamma_\infty \cdot S = \mathrm{PSp}_4(\mathbb{Z})$}

Slide-840 reads off: $A \cdot B$ transforms as a Jacobi-group modular form with a
specified character $\nu_\infty : \Gamma_\infty / \{\pm I\} \to \{\pm 1\}$; testing
against the symplectic involution
\[
  S \;=\; \begin{pmatrix} 0 & I_2 \\ I_2 & 0 \end{pmatrix}
\]
yields anti-invariance of the product --- exactly the $\Dfive$-character.
Since $\Gamma_\infty$ and $S$ generate $\mathrm{PSp}_4(\mathbb{Z})$, the product is
automorphic for the full modular group. The slide cuts at 852 before the
weight-5-character conclusion is stated; lines 853--854 close it.

\subsection*{Cycle~2: 8 Gritsenko--Cl\'ery forms?}

\textbf{Not enumerated in this block.} Gritsenko--Cl\'ery 2008
(arXiv:0812.3962) exhibits exactly $8$ paramodular forms of weight~$5$
on diagonal-divisor types with even character; this interval only
handles the $N=1$ diagonal $\Dfive$. The remaining $7$ forms (levels
$N \in \{2, 3, 4, 6, 7, 8, \ldots\}$ in the Gritsenko--Cl\'ery
classification) enter the Conjecture~$1$ block beyond line 852.

\subsection*{Cycle~3: bridges to CY$_3$}

No explicit CY$_3$ bridge is stated in 723--852. The implicit bridge
--- present in the Fourier coefficient $f(0, 0) = 10$, the weight $5$,
and the $1/64$ prefactor --- is Oberdieck--Pixton 2017: the
refined/motivic DT zeta of $K3 \times E$ is $C/\Dfive^2$. This is the
$\kch(K3 \times E) = 0$ (Calabi--Yau, K\"unneth-multiplicative) versus
$\kBKM(\Phi_1) = 5$ (Borcherds weight) non-coincidence captured in the
universal identity. The bridge is \emph{data} in this block, not theorem.

\subsection*{Cycle~4: Conjecture~1 setup}

Not present in 723--852. Conjecture~1 requires (i) the full 8-form list
and (ii) the $g_N$-twisted K3 elliptic-genus coefficients as root
multiplicities --- both live after line 852.

\subsection*{Cycle~5: Vol~III frontier hiding here}

Three frontier extractions not yet folded into Waves 11--13:

(a) \textbf{$\eta^9 \nu_{1,1}$ as the $\gDfive$ real-root generator.} The
factor $A$ is the character of the $\mathrm{rk}\,3$ real-root lattice
$\mathbb{Z}\delta_1 \oplus \mathbb{Z}\delta_2 \oplus \mathbb{Z}\delta_3$;
the $\eta^9$ is the $9 = \dim \mathfrak{sl}_3 = \dim \mathrm{adjoint}(\mathrm{SL}_3)$
connection, which suggests a chiral-side identification of the real-root
sector of $\gDfive$ with a level-$1$ $\widehat{\mathfrak{sl}}_3$ piece modulo
$\nu_{1,1}$ corrections. This is the $\kBKM$-side decomposition the
universal theorem promises but does not yet spell out.

(b) \textbf{The $\log B$ Hecke-minus expression as a CoHA shadow.} The
$\sum_m m^{-2}\,\phizo \mid_0 T_-(m)$ structure is the exponentiated
trace of a Hecke-operator action; under $\Phi$ at $d = 2$, this is the
$K3$-cohomological-Hall-algebra tower. The $1/m^2$ weighting is the
Schiffmann--Vasserot shuffle normalisation; confirming this
identification would give a first-principles chiral-side construction
of $\log(B)$.

(c) \textbf{The $f(0, -1) = 1$ odd imaginary simple root.} This single
coefficient is responsible for $\gDfive$ being a \emph{super}algebra
rather than an ordinary BKM. It is the $\mathbb{Z}/2$-grading input to
the K3 Yangian (\ref{part:k3-yangian}) and the super-Yangian
$Y(\mathfrak{gl}(4 \mid 20))$ conjecture; every other Gritsenko--Cl\'ery form $\Phi_N$ with $N > 1$
may or may not carry this extra odd simple root --- a case-by-case
computation is \emph{required} to read off the super/ordinary
dichotomy. This is the Wave~14 frontier: for each $N \in \{2,3,4,6\}$,
determine whether $\mathfrak{g}_{\Phi_N}$ is a superalgebra or a
plain BKM, and whether the odd-root count matches the
$g_N$-twisted K3 elliptic genus singular coefficient.
